Use \wikilink{s18--b010}{the classification theorem},
\wikilink{s18--b012}{its uniqueness statement}, and
\wikilink{s19--b007}{primary decomposition}.

\begin{proposition}\label{primary-count-invariants}
For a finitely generated abelian group $A$, its free rank together
with the numbers $|A[p^k]|$, for all primes $p$ and $k\ge1$, determines
its isomorphism class. Here $A[p^k]=\{a:p^ka=0\}$ is finite.
\end{proposition}
\begin{proof}\label{primary-count-invariants-proof}
Decompose the torsion part into prime-power cyclic groups using
classification and the Chinese remainder decomposition. For the
$p$-part, write the exponents as $e_1,\ldots,e_r>0$. Then
\[
 c_p(k):=\log_p|A[p^k]|=\sum_i\min(k,e_i),\qquad c_p(0)=0.
\]
The difference $c_p(k)-c_p(k-1)$ counts the exponents at least $k$.
Successive differences therefore recover the multiplicity of every
exponent. The free part contributes no elements to $A[p^k]$, and its
rank is specified separately. Classification now determines $A$.
\end{proof}
\begin{corollary}\label{abelian-cancellation}
If $G,H,K$ are finitely generated abelian groups and
$G\times H\cong G\times K$, then $H\cong K$.
\end{corollary}
\begin{proof}\label{abelian-cancellation-proof}
Ranks add under finite products, so the ranks of $H$ and $K$ agree.
Taking the subgroups annihilated by $p^k$ in the isomorphism and
comparing their finite orders gives
$|G[p^k]|\,|H[p^k]|=|G[p^k]|\,|K[p^k]|$.
Cancel the positive finite factor and apply the proposition.
\end{proof}
\begin{remark}\label{cancellation-scope}
Finite generation is essential: for $F=\bigoplus_{n\ge1}\mathbb Z$,
reindexing a basis gives $F\oplus\mathbb Z\cong F$, although
$\mathbb Z$ is not the zero group. Also, concatenating two
invariant-factor lists need not leave them in divisibility order;
the prime-power invariants avoid this issue.
\end{remark}
